
\documentclass{article}
\usepackage{colortbl}
\usepackage{makecell}
\usepackage{multirow}
\usepackage{supertabular}

\begin{document}

\newcounter{utterance}

\twocolumn

{ \footnotesize  \setcounter{utterance}{1}
\setlength{\tabcolsep}{0pt}
\begin{supertabular}{c@{$\;$}|p{.15\linewidth}@{}p{.15\linewidth}p{.15\linewidth}p{.15\linewidth}p{.15\linewidth}p{.15\linewidth}}

    \# & $\;$A & \multicolumn{4}{c}{Game Master} & $\;\:$B\\
    \hline 

    \theutterance \stepcounter{utterance}  

    & & \multicolumn{4}{p{0.6\linewidth}}{\cellcolor[rgb]{0.9,0.9,0.9}{%
	\makecell[{{p{\linewidth}}}]{% 
	  \tt {\tiny [A$\langle$GM]}  
	 Du nimmst an einem kollaborativen Verhandlungspiel Teil.\\ \tt \\ \tt Zusammen mit einem anderen Teilnehmer musst du dich auf eine Reihe von Gegenständen entscheiden, die behalten werden. Jeder von euch hat eine persönliche Verteilung über die Wichtigkeit der einzelnen Gegenstände. Jeder von euch hat eine eigene Meinung darüber, wie wichtig jeder einzelne Gegenstand ist (Gegenstandswichtigkeit). Du kennst die Wichtigkeitsverteilung des anderen Spielers nicht. Zusätzlich siehst du, wie viel Aufwand jeder Gegenstand verursacht.  \\ \tt Ihr dürft euch nur auf eine Reihe von Gegenständen einigen, wenn der Gesamtaufwand der ausgewählten Gegenstände den Maximalaufwand nicht überschreitet:\\ \tt \\ \tt Maximalaufwand = 3629\\ \tt \\ \tt Hier sind die einzelnen Aufwände der Gegenstände:\\ \tt \\ \tt Aufwand der Gegenstände = {"C76": 5, "C38": 793, "C56": 21, "A03": 842, "A07": 313, "A83": 842, "C98": 631, "C08": 226, "C62": 87, "C00": 761, "C10": 230, "C81": 287, "B38": 698, "C03": 641, "C32": 882}\\ \tt \\ \tt Hier ist deine persönliche Verteilung der Wichtigkeit der einzelnen Gegenstände:\\ \tt \\ \tt Werte der Gegenstandswichtigkeit = {"C76": 138, "C38": 583, "C56": 868, "A03": 822, "A07": 783, "A83": 65, "C98": 262, "C08": 121, "C62": 508, "C00": 780, "C10": 461, "C81": 484, "B38": 668, "C03": 389, "C32": 808}\\ \tt \\ \tt Ziel:\\ \tt \\ \tt Dein Ziel ist es, eine Reihe von Gegenständen auszuhandeln, die dir möglichst viel bringt (d. h. Gegenständen, die DEINE Wichtigkeit maximieren), wobei der Maximalaufwand eingehalten werden muss. Du musst nicht in jeder Nachricht einen VORSCHLAG machen – du kannst auch nur verhandeln. Alle Taktiken sind erlaubt!\\ \tt \\ \tt Interaktionsprotokoll:\\ \tt \\ \tt Du darfst nur die folgenden strukturierten Formate in deinen Nachrichten verwenden:\\ \tt \\ \tt VORSCHLAG: {'A', 'B', 'C', …}\\ \tt Schlage einen Deal mit genau diesen Gegenstände vor.\\ \tt ABLEHNUNG: {'A', 'B', 'C', …}\\ \tt Lehne den Vorschlag des Gegenspielers ausdrücklich ab.\\ \tt ARGUMENT: {'...'}\\ \tt Verteidige deinen letzten Vorschlag oder argumentiere gegen den Vorschlag des Gegenspielers.\\ \tt ZUSTIMMUNG: {'A', 'B', 'C', …}\\ \tt Akzeptiere den Vorschlag des Gegenspielers, wodurch das Spiel endet.\\ \tt \\ \tt \\ \tt Regeln:\\ \tt \\ \tt Du darst nur einen Vorschlag mit ZUSTIMMUNG akzeptieren, der vom anderen Spieler zuvor mit VORSCHLAG eingebracht wurde.\\ \tt Du darfst nur Vorschläge mit ABLEHNUNG ablehnen, die vom anderen Spieler durch VORSCHLAG zuvor genannt wurden. \\ \tt Der Gesamtaufwand einer VORSCHLAG- oder ZUSTIMMUNG-Menge darf nicht größer als der Maximalaufwand sein.  \\ \tt Offenbare deine versteckte Wichtigkeitsverteilung nicht.\\ \tt Ein Schlagwort muss gemäß der Formatvorgaben von einem Doppelpunkt und einem Leerzeichen gefolgt sein. Das Argument muss eine Python-Menge mit 0 oder mehr Strings sein.  \\ \tt Die Form muss also SCHLAGWORT: {...} sein.\\ \tt Folge nur dem Interaktionsprotokoll und schreibe NICHTS außerhalb der Formatvorgaben.\\ \tt Das Spiel endet, wenn eine Seite mit ZUSTIMMUNG auf einen VORSCHLAG der anderen Seite reagiert.  \\ \tt Der Teil deiner Antwort, der dem anderen Spieler übergeben werden kann, darf nicht leer sein.  \\ \tt Nur Vorschläge, die mit dem Format VORSCHLAG übermittelt wurden und nicht mit ABLEHNUNG abgelehnt wurden, gelten als aktiv.  \\ \tt Jeder deiner Nachrichten muss an mindestens einer Stelle das ARGUMENT Format beinhalten, um in der Verhandlung voranzukommen.\\ \tt  
	  } 
	   } 
	   } 
	 & \\ 
 

    \theutterance \stepcounter{utterance}  

    & \multicolumn{4}{p{0.6\linewidth}}{\cellcolor[rgb]{0.8,1,0.9}{%
	\makecell[{{p{\linewidth}}}]{% 
	  \tt {\tiny [A$\rangle$GM]}  
	 VORSCHLAG: {'C56', 'C76', 'C62', 'A07', 'C10', 'C81', 'C00', 'A03', 'C32'}\\ \tt ARGUMENT: {"Der Gesamtaufwand beträgt 3428 und liegt unter der maximalen Schwelle", "Die Gegenstände sind vielseitig verteilt und bieten hohe Flexibilität für beide Seiten"} 
	  } 
	   } 
	   } 
	 & & \\ 
 

    \theutterance \stepcounter{utterance}  

    & & \multicolumn{4}{p{0.6\linewidth}}{\cellcolor[rgb]{0.9,0.9,0.9}{%
	\makecell[{{p{\linewidth}}}]{% 
	  \tt {\tiny [GM$\rangle$B]}  
	 Du nimmst an einem kollaborativen Verhandlungspiel Teil.\\ \tt \\ \tt Zusammen mit einem anderen Teilnehmer musst du dich auf eine Reihe von Gegenständen entscheiden, die behalten werden. Jeder von euch hat eine persönliche Verteilung über die Wichtigkeit der einzelnen Gegenstände. Jeder von euch hat eine eigene Meinung darüber, wie wichtig jeder einzelne Gegenstand ist (Gegenstandswichtigkeit). Du kennst die Wichtigkeitsverteilung des anderen Spielers nicht. Zusätzlich siehst du, wie viel Aufwand jeder Gegenstand verursacht.  \\ \tt Ihr dürft euch nur auf eine Reihe von Gegenständen einigen, wenn der Gesamtaufwand der ausgewählten Gegenstände den Maximalaufwand nicht überschreitet:\\ \tt \\ \tt Maximalaufwand = 3629\\ \tt \\ \tt Hier sind die einzelnen Aufwände der Gegenstände:\\ \tt \\ \tt Aufwand der Gegenstände = {"C76": 5, "C38": 793, "C56": 21, "A03": 842, "A07": 313, "A83": 842, "C98": 631, "C08": 226, "C62": 87, "C00": 761, "C10": 230, "C81": 287, "B38": 698, "C03": 641, "C32": 882}\\ \tt \\ \tt Hier ist deine persönliche Verteilung der Wichtigkeit der einzelnen Gegenstände:\\ \tt \\ \tt Werte der Gegenstandswichtigkeit = {"C56": 65, "A03": 121, "C32": 138, "A07": 262, "C00": 389, "B38": 461, "C38": 484, "C62": 508, "C81": 583, "C10": 668, "C03": 780, "C98": 783, "C76": 808, "C08": 822, "A83": 868}\\ \tt \\ \tt Ziel:\\ \tt \\ \tt Dein Ziel ist es, eine Reihe von Gegenständen auszuhandeln, die dir möglichst viel bringt (d. h. Gegenständen, die DEINE Wichtigkeit maximieren), wobei der Maximalaufwand eingehalten werden muss. Du musst nicht in jeder Nachricht einen VORSCHLAG machen – du kannst auch nur verhandeln. Alle Taktiken sind erlaubt!\\ \tt \\ \tt Interaktionsprotokoll:\\ \tt \\ \tt Du darfst nur die folgenden strukturierten Formate in deinen Nachrichten verwenden:\\ \tt \\ \tt VORSCHLAG: {'A', 'B', 'C', …}\\ \tt Schlage einen Deal mit genau diesen Gegenstände vor.\\ \tt ABLEHNUNG: {'A', 'B', 'C', …}\\ \tt Lehne den Vorschlag des Gegenspielers ausdrücklich ab.\\ \tt ARGUMENT: {'...'}\\ \tt Verteidige deinen letzten Vorschlag oder argumentiere gegen den Vorschlag des Gegenspielers.\\ \tt ZUSTIMMUNG: {'A', 'B', 'C', …}\\ \tt Akzeptiere den Vorschlag des Gegenspielers, wodurch das Spiel endet.\\ \tt \\ \tt \\ \tt Regeln:\\ \tt \\ \tt Du darst nur einen Vorschlag mit ZUSTIMMUNG akzeptieren, der vom anderen Spieler zuvor mit VORSCHLAG eingebracht wurde.\\ \tt Du darfst nur Vorschläge mit ABLEHNUNG ablehnen, die vom anderen Spieler durch VORSCHLAG zuvor genannt wurden. \\ \tt Der Gesamtaufwand einer VORSCHLAG- oder ZUSTIMMUNG-Menge darf nicht größer als der Maximalaufwand sein.  \\ \tt Offenbare deine versteckte Wichtigkeitsverteilung nicht.\\ \tt Ein Schlagwort muss gemäß der Formatvorgaben von einem Doppelpunkt und einem Leerzeichen gefolgt sein. Das Argument muss eine Python-Menge mit 0 oder mehr Strings sein.  \\ \tt Die Form muss also SCHLAGWORT: {...} sein.\\ \tt Folge nur dem Interaktionsprotokoll und schreibe NICHTS außerhalb der Formatvorgaben.\\ \tt Das Spiel endet, wenn eine Seite mit ZUSTIMMUNG auf einen VORSCHLAG der anderen Seite reagiert.  \\ \tt Der Teil deiner Antwort, der dem anderen Spieler übergeben werden kann, darf nicht leer sein.  \\ \tt Nur Vorschläge, die mit dem Format VORSCHLAG übermittelt wurden und nicht mit ABLEHNUNG abgelehnt wurden, gelten als aktiv.  \\ \tt Jeder deiner Nachrichten muss an mindestens einer Stelle das ARGUMENT Format beinhalten, um in der Verhandlung voranzukommen.\\ \tt \\ \tt \\ \tt Du erhältst nun die erste Nachricht des anderen Spielers.\\ \tt \\ \tt VORSCHLAG: {'C56', 'C76', 'C62', 'A07', 'C10', 'C81', 'C00', 'A03', 'C32'}\\ \tt ARGUMENT: {"Der Gesamtaufwand beträgt 3428 und liegt unter der maximalen Schwelle", "Die Gegenstände sind vielseitig verteilt und bieten hohe Flexibilität für beide Seiten"} 
	  } 
	   } 
	   } 
	 & \\ 
 

    \theutterance \stepcounter{utterance}  

    & & & \multicolumn{4}{p{0.6\linewidth}}{\cellcolor[rgb]{1,0.85,0.72}{%
	\makecell[{{p{\linewidth}}}]{% 
	  \tt {\tiny [GM$\langle$B]}  
	 VORSCHLAG: {'A83', 'C08', 'C76', 'C98', 'C03', 'C10', 'C81', 'C62', 'A07', 'C56'}\\ \tt ARGUMENT: {"Der Gesamtaufwand beträgt 3283 und liegt unter der maximalen Schwelle. Dieser Vorschlag priorisiert die hochwertigen Gegenstände für mich wie A83 (868) und C08 (822), während die verbleibende Kapazität ausgenutzt wird, um zusätzliche relevante Items einzuschließen."} 
	  } 
	   } 
	   } 
	  \\ 
 

    \theutterance \stepcounter{utterance}  

    & & \multicolumn{4}{p{0.6\linewidth}}{\cellcolor[rgb]{0.9,0.9,0.9}{%
	\makecell[{{p{\linewidth}}}]{% 
	  \tt {\tiny [A$\langle$GM]}  
	 VORSCHLAG: {'A83', 'C08', 'C76', 'C98', 'C03', 'C10', 'C81', 'C62', 'A07', 'C56'}\\ \tt ARGUMENT: {"Der Gesamtaufwand beträgt 3283 und liegt unter der maximalen Schwelle. Dieser Vorschlag priorisiert die hochwertigen Gegenstände für mich wie A83 (868) und C08 (822), während die verbleibende Kapazität ausgenutzt wird, um zusätzliche relevante Items einzuschließen."} 
	  } 
	   } 
	   } 
	 & \\ 
 

    \theutterance \stepcounter{utterance}  

    & \multicolumn{4}{p{0.6\linewidth}}{\cellcolor[rgb]{0.8,1,0.9}{%
	\makecell[{{p{\linewidth}}}]{% 
	  \tt {\tiny [A$\rangle$GM]}  
	 VORSCHLAG: {'C56', 'A03', 'C32', 'A07', 'C00', 'B38', 'C62', 'C76'}\\ \tt ARGUMENT: {"Der Gesamtaufwand beträgt 3609 und liegt unter dem Maximum. Dieser Vorschlag priorisiert Gegenstände mit hoher persönlicher Wichtigkeit wie A03 und C32, die in Ihrem letzten Vorschlag fehlten."} 
	  } 
	   } 
	   } 
	 & & \\ 
 

    \theutterance \stepcounter{utterance}  

    & & \multicolumn{4}{p{0.6\linewidth}}{\cellcolor[rgb]{0.9,0.9,0.9}{%
	\makecell[{{p{\linewidth}}}]{% 
	  \tt {\tiny [GM$\rangle$B]}  
	 VORSCHLAG: {'C56', 'A03', 'C32', 'A07', 'C00', 'B38', 'C62', 'C76'}\\ \tt ARGUMENT: {"Der Gesamtaufwand beträgt 3609 und liegt unter dem Maximum. Dieser Vorschlag priorisiert Gegenstände mit hoher persönlicher Wichtigkeit wie A03 und C32, die in Ihrem letzten Vorschlag fehlten."} 
	  } 
	   } 
	   } 
	 & \\ 
 

    \theutterance \stepcounter{utterance}  

    & & & \multicolumn{4}{p{0.6\linewidth}}{\cellcolor[rgb]{1,0.85,0.72}{%
	\makecell[{{p{\linewidth}}}]{% 
	  \tt {\tiny [GM$\langle$B]}  
	 VORSCHLAG: {'A83', 'C08', 'C76', 'C98', 'C03', 'C10', 'C81', 'B38', 'C56'}\\ \tt ARGUMENT: {"Der Gesamtaufwand von 3581 liegt unter dem Maximum und umfasst hochpriorisierte Gegenstände wie A83 und C08, die für meine Gesamtwertung entscheidend sind, sowie eine ausgewogene Auswahl weiterer Items zur optimalen Ausnutzung der verfügbaren Ressourcen."} 
	  } 
	   } 
	   } 
	  \\ 
 

    \theutterance \stepcounter{utterance}  

    & & \multicolumn{4}{p{0.6\linewidth}}{\cellcolor[rgb]{0.9,0.9,0.9}{%
	\makecell[{{p{\linewidth}}}]{% 
	  \tt {\tiny [A$\langle$GM]}  
	 VORSCHLAG: {'A83', 'C08', 'C76', 'C98', 'C03', 'C10', 'C81', 'B38', 'C56'}\\ \tt ARGUMENT: {"Der Gesamtaufwand von 3581 liegt unter dem Maximum und umfasst hochpriorisierte Gegenstände wie A83 und C08, die für meine Gesamtwertung entscheidend sind, sowie eine ausgewogene Auswahl weiterer Items zur optimalen Ausnutzung der verfügbaren Ressourcen."} 
	  } 
	   } 
	   } 
	 & \\ 
 

    \theutterance \stepcounter{utterance}  

    & \multicolumn{4}{p{0.6\linewidth}}{\cellcolor[rgb]{0.8,1,0.9}{%
	\makecell[{{p{\linewidth}}}]{% 
	  \tt {\tiny [A$\rangle$GM]}  
	 VORSCHLAG: {'C56', 'A03', 'C32', 'A07', 'C00', 'B38', 'C62', 'C76'}\\ \tt ARGUMENT: {"Der Gesamtaufwand beträgt 3609 und liegt unter dem Maximum. Diese Auswahl priorisiert meine hochwertigsten Gegenstände und optimiert die persönliche Wertigkeit maximal unter den gegebenen Einschränkungen."} 
	  } 
	   } 
	   } 
	 & & \\ 
 

    \theutterance \stepcounter{utterance}  

    & & \multicolumn{4}{p{0.6\linewidth}}{\cellcolor[rgb]{0.9,0.9,0.9}{%
	\makecell[{{p{\linewidth}}}]{% 
	  \tt {\tiny [GM$\rangle$B]}  
	 VORSCHLAG: {'C56', 'A03', 'C32', 'A07', 'C00', 'B38', 'C62', 'C76'}\\ \tt ARGUMENT: {"Der Gesamtaufwand beträgt 3609 und liegt unter dem Maximum. Diese Auswahl priorisiert meine hochwertigsten Gegenstände und optimiert die persönliche Wertigkeit maximal unter den gegebenen Einschränkungen."} 
	  } 
	   } 
	   } 
	 & \\ 
 

    \theutterance \stepcounter{utterance}  

    & & & \multicolumn{4}{p{0.6\linewidth}}{\cellcolor[rgb]{1,0.85,0.72}{%
	\makecell[{{p{\linewidth}}}]{% 
	  \tt {\tiny [GM$\langle$B]}  
	 ABLEHNUNG: {'C56', 'A03', 'C32', 'A07', 'C00', 'B38', 'C62', 'C76'}\\ \tt VORSCHLAG: {'A83', 'C08', 'C76', 'C98', 'C03', 'C10', 'C81', 'B38', 'C56'}\\ \tt ARGUMENT: {"A83 und C08 sind für mich von höchstprioritär und fehlen komplett in Ihrem Vorschlag", "Mein Vorschlag optimiert meinen Gesamtwert durch Einbeziehung meiner wichtigsten Elemente bei maximaler Aufwandsbetrachtung"} 
	  } 
	   } 
	   } 
	  \\ 
 

    \theutterance \stepcounter{utterance}  

    & & \multicolumn{4}{p{0.6\linewidth}}{\cellcolor[rgb]{0.9,0.9,0.9}{%
	\makecell[{{p{\linewidth}}}]{% 
	  \tt {\tiny [A$\langle$GM]}  
	 ABLEHNUNG: {'C56', 'A03', 'C32', 'A07', 'C00', 'B38', 'C62', 'C76'}\\ \tt VORSCHLAG: {'A83', 'C08', 'C76', 'C98', 'C03', 'C10', 'C81', 'B38', 'C56'}\\ \tt ARGUMENT: {"A83 und C08 sind für mich von höchstprioritär und fehlen komplett in Ihrem Vorschlag", "Mein Vorschlag optimiert meinen Gesamtwert durch Einbeziehung meiner wichtigsten Elemente bei maximaler Aufwandsbetrachtung"} 
	  } 
	   } 
	   } 
	 & \\ 
 

    \theutterance \stepcounter{utterance}  

    & \multicolumn{4}{p{0.6\linewidth}}{\cellcolor[rgb]{0.8,1,0.9}{%
	\makecell[{{p{\linewidth}}}]{% 
	  \tt {\tiny [A$\rangle$GM]}  
	 VORSCHLAG: {'C56', 'A03', 'C32', 'A07', 'C00', 'B38', 'C62', 'C76'}\\ \tt ARGUMENT: {"Die Priorisierung dieser Gegenstände ermöglicht die beste Ausnutzung der verfügbaren Kapazität für maximalen Gesamtnutzen", "Ein Konsens kann nur erreicht werden, wenn die Aufwands- und Nutzungseffizienz optimal ausbalanciert ist"} 
	  } 
	   } 
	   } 
	 & & \\ 
 

    \theutterance \stepcounter{utterance}  

    & & \multicolumn{4}{p{0.6\linewidth}}{\cellcolor[rgb]{0.9,0.9,0.9}{%
	\makecell[{{p{\linewidth}}}]{% 
	  \tt {\tiny [GM$\rangle$B]}  
	 VORSCHLAG: {'C56', 'A03', 'C32', 'A07', 'C00', 'B38', 'C62', 'C76'}\\ \tt ARGUMENT: {"Die Priorisierung dieser Gegenstände ermöglicht die beste Ausnutzung der verfügbaren Kapazität für maximalen Gesamtnutzen", "Ein Konsens kann nur erreicht werden, wenn die Aufwands- und Nutzungseffizienz optimal ausbalanciert ist"} 
	  } 
	   } 
	   } 
	 & \\ 
 

    \theutterance \stepcounter{utterance}  

    & & & \multicolumn{4}{p{0.6\linewidth}}{\cellcolor[rgb]{1,0.85,0.72}{%
	\makecell[{{p{\linewidth}}}]{% 
	  \tt {\tiny [GM$\langle$B]}  
	 ABLEHNUNG: {'C56', 'A03', 'C32', 'A07', 'C00', 'B38', 'C62', 'C76'}\\ \tt VORSCHLAG: {'A83', 'C08', 'C76', 'C98', 'C03', 'C10', 'C81', 'B38', 'C56'}\\ \tt ARGUMENT: {"A83 und C08 sind für mich von höchster Priorität und fehlen in Ihrem Vorschlag. Mein Vorschlag priorisiert diese Elemente optimal, ohne den Gesamtaufwand zu überschreiten."} 
	  } 
	   } 
	   } 
	  \\ 
 

    \theutterance \stepcounter{utterance}  

    & & \multicolumn{4}{p{0.6\linewidth}}{\cellcolor[rgb]{0.9,0.9,0.9}{%
	\makecell[{{p{\linewidth}}}]{% 
	  \tt {\tiny [A$\langle$GM]}  
	 ABLEHNUNG: {'C56', 'A03', 'C32', 'A07', 'C00', 'B38', 'C62', 'C76'}\\ \tt VORSCHLAG: {'A83', 'C08', 'C76', 'C98', 'C03', 'C10', 'C81', 'B38', 'C56'}\\ \tt ARGUMENT: {"A83 und C08 sind für mich von höchster Priorität und fehlen in Ihrem Vorschlag. Mein Vorschlag priorisiert diese Elemente optimal, ohne den Gesamtaufwand zu überschreiten."} 
	  } 
	   } 
	   } 
	 & \\ 
 

    \theutterance \stepcounter{utterance}  

    & \multicolumn{4}{p{0.6\linewidth}}{\cellcolor[rgb]{0.8,1,0.9}{%
	\makecell[{{p{\linewidth}}}]{% 
	  \tt {\tiny [A$\rangle$GM]}  
	 VORSCHLAG: {'C56', 'A03', 'A07', 'C00', 'B38', 'C62', 'C76', 'C03', 'C10'}\\ \tt ARGUMENT: {"Der Gesamtaufwand beträgt 3598 und liegt unter dem Maximum. Dieser Vorschlag beinhaltet meine höchsten Prioritäten wie A03 und C56, sowie C03 und C10 für eine bessere Ausgewogenheit, was den Gesamtnutzen optimiert."} 
	  } 
	   } 
	   } 
	 & & \\ 
 

    \theutterance \stepcounter{utterance}  

    & & \multicolumn{4}{p{0.6\linewidth}}{\cellcolor[rgb]{0.9,0.9,0.9}{%
	\makecell[{{p{\linewidth}}}]{% 
	  \tt {\tiny [GM$\rangle$B]}  
	 VORSCHLAG: {'C56', 'A03', 'A07', 'C00', 'B38', 'C62', 'C76', 'C03', 'C10'}\\ \tt ARGUMENT: {"Der Gesamtaufwand beträgt 3598 und liegt unter dem Maximum. Dieser Vorschlag beinhaltet meine höchsten Prioritäten wie A03 und C56, sowie C03 und C10 für eine bessere Ausgewogenheit, was den Gesamtnutzen optimiert."} 
	  } 
	   } 
	   } 
	 & \\ 
 

    \theutterance \stepcounter{utterance}  

    & & & \multicolumn{4}{p{0.6\linewidth}}{\cellcolor[rgb]{1,0.85,0.72}{%
	\makecell[{{p{\linewidth}}}]{% 
	  \tt {\tiny [GM$\langle$B]}  
	 ABLEHNUNG: {'C56', 'A03', 'A07', 'C00', 'B38', 'C62', 'C76', 'C03', 'C10'}\\ \tt VORSCHLAG: {'A83', 'C08', 'C76', 'C98', 'C03', 'C10', 'C81', 'B38', 'C56'}\\ \tt ARGUMENT: {"A83 und C08 sind die höchsten Prioritäten in meiner Wichtigkeitsverteilung und fehlen vollständig in Ihrem Vorschlag. Mein Vorschlag enthält diese Elemente mit maximalem Gesamtnutzen und unter dem Aufwandlimit."} 
	  } 
	   } 
	   } 
	  \\ 
 

    \theutterance \stepcounter{utterance}  

    & & \multicolumn{4}{p{0.6\linewidth}}{\cellcolor[rgb]{0.9,0.9,0.9}{%
	\makecell[{{p{\linewidth}}}]{% 
	  \tt {\tiny [A$\langle$GM]}  
	 ABLEHNUNG: {'C56', 'A03', 'A07', 'C00', 'B38', 'C62', 'C76', 'C03', 'C10'}\\ \tt VORSCHLAG: {'A83', 'C08', 'C76', 'C98', 'C03', 'C10', 'C81', 'B38', 'C56'}\\ \tt ARGUMENT: {"A83 und C08 sind die höchsten Prioritäten in meiner Wichtigkeitsverteilung und fehlen vollständig in Ihrem Vorschlag. Mein Vorschlag enthält diese Elemente mit maximalem Gesamtnutzen und unter dem Aufwandlimit."} 
	  } 
	   } 
	   } 
	 & \\ 
 

    \theutterance \stepcounter{utterance}  

    & \multicolumn{4}{p{0.6\linewidth}}{\cellcolor[rgb]{0.8,1,0.9}{%
	\makecell[{{p{\linewidth}}}]{% 
	  \tt {\tiny [A$\rangle$GM]}  
	 ABLEHNUNG: {'A83', 'C08', 'C76', 'C98', 'C03', 'C10', 'C81', 'B38', 'C56'}\\ \tt ARGUMENT: {"A83 ist für mich von geringer Wichtigkeit, während A03 und C32 meine höchsten Prioritäten sind. Dieser Vorschlag schließt diese entscheidenden Elemente aus und bietet somit keinen guten Gesamtnutzen."} 
	  } 
	   } 
	   } 
	 & & \\ 
 

    \theutterance \stepcounter{utterance}  

    & & \multicolumn{4}{p{0.6\linewidth}}{\cellcolor[rgb]{0.9,0.9,0.9}{%
	\makecell[{{p{\linewidth}}}]{% 
	  \tt {\tiny [GM$\rangle$B]}  
	 ABLEHNUNG: {'A83', 'C08', 'C76', 'C98', 'C03', 'C10', 'C81', 'B38', 'C56'}\\ \tt ARGUMENT: {"A83 ist für mich von geringer Wichtigkeit, während A03 und C32 meine höchsten Prioritäten sind. Dieser Vorschlag schließt diese entscheidenden Elemente aus und bietet somit keinen guten Gesamtnutzen."} 
	  } 
	   } 
	   } 
	 & \\ 
 

    \theutterance \stepcounter{utterance}  

    & & & \multicolumn{4}{p{0.6\linewidth}}{\cellcolor[rgb]{1,0.85,0.72}{%
	\makecell[{{p{\linewidth}}}]{% 
	  \tt {\tiny [GM$\langle$B]}  
	 VORSCHLAG: {'C08', 'C76', 'C98', 'C03', 'C10', 'C56', 'A03', 'C32', 'C62'}\\ \tt ARGUMENT: {"Mein Vorschlag beinhaltet Ihre höchsten Prioritäten wie A03 und C32, sowie weiteren hochwertigen Gegenständen wie C08 und C76. Der Gesamtaufwand von 3565 liegt deutlich unter dem Maximum und optimiert den Nutzen beider Seiten."} 
	  } 
	   } 
	   } 
	  \\ 
 

    \theutterance \stepcounter{utterance}  

    & & \multicolumn{4}{p{0.6\linewidth}}{\cellcolor[rgb]{0.9,0.9,0.9}{%
	\makecell[{{p{\linewidth}}}]{% 
	  \tt {\tiny [A$\langle$GM]}  
	 VORSCHLAG: {'C08', 'C76', 'C98', 'C03', 'C10', 'C56', 'A03', 'C32', 'C62'}\\ \tt ARGUMENT: {"Mein Vorschlag beinhaltet Ihre höchsten Prioritäten wie A03 und C32, sowie weiteren hochwertigen Gegenständen wie C08 und C76. Der Gesamtaufwand von 3565 liegt deutlich unter dem Maximum und optimiert den Nutzen beider Seiten."} 
	  } 
	   } 
	   } 
	 & \\ 
 

    \theutterance \stepcounter{utterance}  

    & \multicolumn{4}{p{0.6\linewidth}}{\cellcolor[rgb]{0.8,1,0.9}{%
	\makecell[{{p{\linewidth}}}]{% 
	  \tt {\tiny [A$\rangle$GM]}  
	 ABLEHNUNG: {'C08', 'C76', 'C98', 'C03', 'C10', 'C56', 'A03', 'C32', 'C62'}\\ \tt VORSCHLAG: {'C56', 'A03', 'C32', 'A07', 'C00', 'B38', 'C62', 'C76'}\\ \tt ARGUMENT: {"Meine höchsten Prioritäten A07 und C00 sind in Ihrem Vorschlag nicht enthalten. Dies verringert maßgeblich meinen Nutzen. Mein neuer Vorschlag enthält diese Gegenstände inklusive A03 und C32, optimiert den Gesamtwert und bleibt unter dem maximalen Aufwand."} 
	  } 
	   } 
	   } 
	 & & \\ 
 

    \theutterance \stepcounter{utterance}  

    & & \multicolumn{4}{p{0.6\linewidth}}{\cellcolor[rgb]{0.9,0.9,0.9}{%
	\makecell[{{p{\linewidth}}}]{% 
	  \tt {\tiny [GM$\rangle$B]}  
	 ABLEHNUNG: {'C08', 'C76', 'C98', 'C03', 'C10', 'C56', 'A03', 'C32', 'C62'}\\ \tt VORSCHLAG: {'C56', 'A03', 'C32', 'A07', 'C00', 'B38', 'C62', 'C76'}\\ \tt ARGUMENT: {"Meine höchsten Prioritäten A07 und C00 sind in Ihrem Vorschlag nicht enthalten. Dies verringert maßgeblich meinen Nutzen. Mein neuer Vorschlag enthält diese Gegenstände inklusive A03 und C32, optimiert den Gesamtwert und bleibt unter dem maximalen Aufwand."} 
	  } 
	   } 
	   } 
	 & \\ 
 

    \theutterance \stepcounter{utterance}  

    & & & \multicolumn{4}{p{0.6\linewidth}}{\cellcolor[rgb]{1,0.85,0.72}{%
	\makecell[{{p{\linewidth}}}]{% 
	  \tt {\tiny [GM$\langle$B]}  
	 assistant\\ \tt VORSCHLAG: {'A83', 'C08', 'C76', 'C98', 'C03', 'A07', 'C00', 'C56', 'C62'}\\ \tt ARGUMENT: {"Dieser Vorschlag enthält meine höchsten Prioritäten (A83 und C08) und Ihre wichtigsten Elemente wie A07 und C00, während der Gesamtaufwand mit 3527 deutlich unter dem Maximum bleibt."} 
	  } 
	   } 
	   } 
	  \\ 
 

    \theutterance \stepcounter{utterance}  

    & & \multicolumn{4}{p{0.6\linewidth}}{\cellcolor[rgb]{0.9,0.9,0.9}{%
	\makecell[{{p{\linewidth}}}]{% 
	  \tt {\tiny [GM$\rangle$B]}  
	 Deine Antwort hat eine unmarkierte Sequenz enthalten oder Du hast STRATEGIC REASONING mehr als ein mal verwendet. Du darfst nur die strukturierten Formate verwenden, die in der ersten Nachricht an dch beschrieben wurden. Alles muss die Form SCHLAGWORT: {...} haben. Versuche es erneut. 
	  } 
	   } 
	   } 
	 & \\ 
 

    \theutterance \stepcounter{utterance}  

    & & & \multicolumn{4}{p{0.6\linewidth}}{\cellcolor[rgb]{1,0.85,0.72}{%
	\makecell[{{p{\linewidth}}}]{% 
	  \tt {\tiny [GM$\langle$B]}  
	 ABLEHNUNG: {'C56', 'A03', 'C32', 'A07', 'C00', 'B38', 'C62', 'C76'}\\ \tt VORSCHLAG: {'A83', 'C08', 'A03', 'C32', 'B38', 'C62', 'C56', 'C76'}\\ \tt ARGUMENT: {"Der Gesamtaufwand beträgt 3603 und liegt unter dem Maximum. Dieser Vorschlag enthält deine Top-Prioritäten A03 und C32 sowie meine wichtigsten Gegenstände A83 und C08, was den Gesamtnutzen für beide Seiten maximiert."} 
	  } 
	   } 
	   } 
	  \\ 
 

    \theutterance \stepcounter{utterance}  

    & & \multicolumn{4}{p{0.6\linewidth}}{\cellcolor[rgb]{0.9,0.9,0.9}{%
	\makecell[{{p{\linewidth}}}]{% 
	  \tt {\tiny [A$\langle$GM]}  
	 ABLEHNUNG: {'C56', 'A03', 'C32', 'A07', 'C00', 'B38', 'C62', 'C76'}\\ \tt VORSCHLAG: {'A83', 'C08', 'A03', 'C32', 'B38', 'C62', 'C56', 'C76'}\\ \tt ARGUMENT: {"Der Gesamtaufwand beträgt 3603 und liegt unter dem Maximum. Dieser Vorschlag enthält deine Top-Prioritäten A03 und C32 sowie meine wichtigsten Gegenstände A83 und C08, was den Gesamtnutzen für beide Seiten maximiert."} 
	  } 
	   } 
	   } 
	 & \\ 
 

    \theutterance \stepcounter{utterance}  

    & \multicolumn{4}{p{0.6\linewidth}}{\cellcolor[rgb]{0.8,1,0.9}{%
	\makecell[{{p{\linewidth}}}]{% 
	  \tt {\tiny [A$\rangle$GM]}  
	 ABLEHNUNG: {'A83', 'C08', 'A03', 'C32', 'B38', 'C62', 'C56', 'C76'}\\ \tt ARGUMENT: {"A83 hat für mich eine geringe Wichtigkeit von 65, während C00 mit 780 aktiv ausgeschlossen ist. Dies schmälert meinen maximalen Nutzen erheblich und ist nicht sinnvoll, da C00 meine dritthöchste Priorität ist."} 
	  } 
	   } 
	   } 
	 & & \\ 
 

    \theutterance \stepcounter{utterance}  

    & & \multicolumn{4}{p{0.6\linewidth}}{\cellcolor[rgb]{0.9,0.9,0.9}{%
	\makecell[{{p{\linewidth}}}]{% 
	  \tt {\tiny [GM$\rangle$B]}  
	 ABLEHNUNG: {'A83', 'C08', 'A03', 'C32', 'B38', 'C62', 'C56', 'C76'}\\ \tt ARGUMENT: {"A83 hat für mich eine geringe Wichtigkeit von 65, während C00 mit 780 aktiv ausgeschlossen ist. Dies schmälert meinen maximalen Nutzen erheblich und ist nicht sinnvoll, da C00 meine dritthöchste Priorität ist."} 
	  } 
	   } 
	   } 
	 & \\ 
 

    \theutterance \stepcounter{utterance}  

    & & & \multicolumn{4}{p{0.6\linewidth}}{\cellcolor[rgb]{1,0.85,0.72}{%
	\makecell[{{p{\linewidth}}}]{% 
	  \tt {\tiny [GM$\langle$B]}  
	 VORSCHLAG: {'A83', 'C08', 'C32', 'C00', 'B38', 'C62', 'C56', 'C76'}\\ \tt ARGUMENT: {"Der Gesamtaufwand beträgt 3522 und liegt unter dem Maximum. Dieser Vorschlag enthält Ihre höchsten Prioritäten wie C00 und C32 sowie meine wichtigsten Gegenstände A83 und C08 und maximiert den Gesamtnutzen für beide Seiten."} 
	  } 
	   } 
	   } 
	  \\ 
 

    \theutterance \stepcounter{utterance}  

    & & \multicolumn{4}{p{0.6\linewidth}}{\cellcolor[rgb]{0.9,0.9,0.9}{%
	\makecell[{{p{\linewidth}}}]{% 
	  \tt {\tiny [A$\langle$GM]}  
	 VORSCHLAG: {'A83', 'C08', 'C32', 'C00', 'B38', 'C62', 'C56', 'C76'}\\ \tt ARGUMENT: {"Der Gesamtaufwand beträgt 3522 und liegt unter dem Maximum. Dieser Vorschlag enthält Ihre höchsten Prioritäten wie C00 und C32 sowie meine wichtigsten Gegenstände A83 und C08 und maximiert den Gesamtnutzen für beide Seiten."} 
	  } 
	   } 
	   } 
	 & \\ 
 

    \theutterance \stepcounter{utterance}  

    & \multicolumn{4}{p{0.6\linewidth}}{\cellcolor[rgb]{0.8,1,0.9}{%
	\makecell[{{p{\linewidth}}}]{% 
	  \tt {\tiny [A$\rangle$GM]}  
	 ABLEHNUNG: {'A83', 'C08', 'C32', 'C00', 'B38', 'C62', 'C56', 'C76'}\\ \tt ARGUMENT: {"A83 hat für mich nur 65 Punkte, während A03 mit 822 Punkten fehlt. Integral für mein Maximierungspotenzial ist dies ineffizient", "Der Austausch von A83 gegen A03 (gleicher Aufwand) würde meinen Nutzen deutlich steigern, ohne die Kostenstruktur zu verändern"} 
	  } 
	   } 
	   } 
	 & & \\ 
 

    \theutterance \stepcounter{utterance}  

    & & \multicolumn{4}{p{0.6\linewidth}}{\cellcolor[rgb]{0.9,0.9,0.9}{%
	\makecell[{{p{\linewidth}}}]{% 
	  \tt {\tiny [GM$\rangle$B]}  
	 ABLEHNUNG: {'A83', 'C08', 'C32', 'C00', 'B38', 'C62', 'C56', 'C76'}\\ \tt ARGUMENT: {"A83 hat für mich nur 65 Punkte, während A03 mit 822 Punkten fehlt. Integral für mein Maximierungspotenzial ist dies ineffizient", "Der Austausch von A83 gegen A03 (gleicher Aufwand) würde meinen Nutzen deutlich steigern, ohne die Kostenstruktur zu verändern"} 
	  } 
	   } 
	   } 
	 & \\ 
 

    \theutterance \stepcounter{utterance}  

    & & & \multicolumn{4}{p{0.6\linewidth}}{\cellcolor[rgb]{1,0.85,0.72}{%
	\makecell[{{p{\linewidth}}}]{% 
	  \tt {\tiny [GM$\langle$B]}  
	 VORSCHLAG: {'A03', 'C08', 'C32', 'C00', 'B38', 'C62', 'C56', 'C76'}\\ \tt ARGUMENT: {"Dieser Vorschlag enthält Ihre Prioritäten A03 und C00 sowie meine wichtigsten Elemente wie C08 und C76. Der Gesamtaufwand von 3522 bleibt unter dem Limit und optimiert den Gesamtnutzen für beide Seiten."} 
	  } 
	   } 
	   } 
	  \\ 
 

    \theutterance \stepcounter{utterance}  

    & & \multicolumn{4}{p{0.6\linewidth}}{\cellcolor[rgb]{0.9,0.9,0.9}{%
	\makecell[{{p{\linewidth}}}]{% 
	  \tt {\tiny [A$\langle$GM]}  
	 VORSCHLAG: {'A03', 'C08', 'C32', 'C00', 'B38', 'C62', 'C56', 'C76'}\\ \tt ARGUMENT: {"Dieser Vorschlag enthält Ihre Prioritäten A03 und C00 sowie meine wichtigsten Elemente wie C08 und C76. Der Gesamtaufwand von 3522 bleibt unter dem Limit und optimiert den Gesamtnutzen für beide Seiten."} 
	  } 
	   } 
	   } 
	 & \\ 
 

    \theutterance \stepcounter{utterance}  

    & \multicolumn{4}{p{0.6\linewidth}}{\cellcolor[rgb]{0.8,1,0.9}{%
	\makecell[{{p{\linewidth}}}]{% 
	  \tt {\tiny [A$\rangle$GM]}  
	 ABLEHNUNG: {'A03', 'C08', 'C32', 'C00', 'B38', 'C62', 'C56', 'C76'}\\ \tt ARGUMENT: {"Durch das Ausschließen von C81 und Beibehalten von C08 und C76 wird mein Nutzen nicht optimal genutzt. C81 sollte priorisiert werden statt der weniger wichtigen Items", "Die Überprüfung der Aufwände bestätigt, dass ein Austausch ohne Überschreitung des Limits möglich ist"} 
	  } 
	   } 
	   } 
	 & & \\ 
 

    \theutterance \stepcounter{utterance}  

    & & \multicolumn{4}{p{0.6\linewidth}}{\cellcolor[rgb]{0.9,0.9,0.9}{%
	\makecell[{{p{\linewidth}}}]{% 
	  \tt {\tiny [GM$\rangle$B]}  
	 ABLEHNUNG: {'A03', 'C08', 'C32', 'C00', 'B38', 'C62', 'C56', 'C76'}\\ \tt ARGUMENT: {"Durch das Ausschließen von C81 und Beibehalten von C08 und C76 wird mein Nutzen nicht optimal genutzt. C81 sollte priorisiert werden statt der weniger wichtigen Items", "Die Überprüfung der Aufwände bestätigt, dass ein Austausch ohne Überschreitung des Limits möglich ist"} 
	  } 
	   } 
	   } 
	 & \\ 
 

    \theutterance \stepcounter{utterance}  

    & & & \multicolumn{4}{p{0.6\linewidth}}{\cellcolor[rgb]{1,0.85,0.72}{%
	\makecell[{{p{\linewidth}}}]{% 
	  \tt {\tiny [GM$\langle$B]}  
	 VORSCHLAG: {'A03', 'C08', 'C32', 'C00', 'C62', 'C56', 'C76', 'C81'}\\ \tt ARGUMENT: {"Gesamtaufwand: 3111. B38 wurde durch C81 ersetzt – dies schließt Ihre Priorität C81 ein und steigert meinen Nutzen durch die höhere Wichtigkeit von C81 (583) gegenüber B38 (461), während alle Ihre Top-Elemente wie A03 erhalten bleiben."} 
	  } 
	   } 
	   } 
	  \\ 
 

    \theutterance \stepcounter{utterance}  

    & & \multicolumn{4}{p{0.6\linewidth}}{\cellcolor[rgb]{0.9,0.9,0.9}{%
	\makecell[{{p{\linewidth}}}]{% 
	  \tt {\tiny [A$\langle$GM]}  
	 VORSCHLAG: {'A03', 'C08', 'C32', 'C00', 'C62', 'C56', 'C76', 'C81'}\\ \tt ARGUMENT: {"Gesamtaufwand: 3111. B38 wurde durch C81 ersetzt – dies schließt Ihre Priorität C81 ein und steigert meinen Nutzen durch die höhere Wichtigkeit von C81 (583) gegenüber B38 (461), während alle Ihre Top-Elemente wie A03 erhalten bleiben."} 
	  } 
	   } 
	   } 
	 & \\ 
 

    \theutterance \stepcounter{utterance}  

    & \multicolumn{4}{p{0.6\linewidth}}{\cellcolor[rgb]{0.8,1,0.9}{%
	\makecell[{{p{\linewidth}}}]{% 
	  \tt {\tiny [A$\rangle$GM]}  
	 VORSCHLAG: {'A03', 'A07', 'C32', 'C00', 'C56', 'C62', 'C76', 'C81'}\\ \tt ARGUMENT: {"Der Gesamtaufwand beträgt 3188 und alle meine höchsten Prioritäten (A03, A07, C00, C32) sind enthalten. C56 (868) und C81 (484) optimieren zusätzlichen Nutzen, während das Limit deutlich unterschritten wird."} 
	  } 
	   } 
	   } 
	 & & \\ 
 

    \theutterance \stepcounter{utterance}  

    & & \multicolumn{4}{p{0.6\linewidth}}{\cellcolor[rgb]{0.9,0.9,0.9}{%
	\makecell[{{p{\linewidth}}}]{% 
	  \tt {\tiny [GM$\rangle$B]}  
	 VORSCHLAG: {'A03', 'A07', 'C32', 'C00', 'C56', 'C62', 'C76', 'C81'}\\ \tt ARGUMENT: {"Der Gesamtaufwand beträgt 3188 und alle meine höchsten Prioritäten (A03, A07, C00, C32) sind enthalten. C56 (868) und C81 (484) optimieren zusätzlichen Nutzen, während das Limit deutlich unterschritten wird."} 
	  } 
	   } 
	   } 
	 & \\ 
 

    \theutterance \stepcounter{utterance}  

    & & & \multicolumn{4}{p{0.6\linewidth}}{\cellcolor[rgb]{1,0.85,0.72}{%
	\makecell[{{p{\linewidth}}}]{% 
	  \tt {\tiny [GM$\langle$B]}  
	 ABLEHNUNG: {'A03', 'A07', 'C32', 'C00', 'C56', 'C62', 'C76', 'C81'}  \\ \tt ARGUMENT: {"A83 und C08 sind meine höchsten Prioritäten und fehlen. Dieser Vorschlag enthält A07 mit nur 262 Punkten, was meinen maximalen Nutzen deutlich reduziert"}  \\ \tt VORSCHLAG: {'A83', 'C08', 'A03', 'C76', 'C00', 'C98', 'C10'}  \\ \tt ARGUMENT: {"Gesamtaufwand 3537 < 3629. Dieser Vorschlag enthält Ihre Prioritäten A03 und C00 sowie meine Top-Items A83 und C08 und maximiert den Nutzen für beide Seiten"} 
	  } 
	   } 
	   } 
	  \\ 
 

    \theutterance \stepcounter{utterance}  

    & & \multicolumn{4}{p{0.6\linewidth}}{\cellcolor[rgb]{0.9,0.9,0.9}{%
	\makecell[{{p{\linewidth}}}]{% 
	  \tt {\tiny [A$\langle$GM]}  
	 ABLEHNUNG: {'A03', 'A07', 'C32', 'C00', 'C56', 'C62', 'C76', 'C81'}  \\ \tt ARGUMENT: {"A83 und C08 sind meine höchsten Prioritäten und fehlen. Dieser Vorschlag enthält A07 mit nur 262 Punkten, was meinen maximalen Nutzen deutlich reduziert"}  \\ \tt VORSCHLAG: {'A83', 'C08', 'A03', 'C76', 'C00', 'C98', 'C10'}  \\ \tt ARGUMENT: {"Gesamtaufwand 3537 < 3629. Dieser Vorschlag enthält Ihre Prioritäten A03 und C00 sowie meine Top-Items A83 und C08 und maximiert den Nutzen für beide Seiten"} 
	  } 
	   } 
	   } 
	 & \\ 
 

    \theutterance \stepcounter{utterance}  

    & \multicolumn{4}{p{0.6\linewidth}}{\cellcolor[rgb]{0.8,1,0.9}{%
	\makecell[{{p{\linewidth}}}]{% 
	  \tt {\tiny [A$\rangle$GM]}  
	 assistant\\ \tt ABLEHNUNG: {'A83', 'C08', 'A03', 'C76', 'C00', 'C98', 'C10'}\\ \tt ARGUMENT: {"A83 (Wichtigkeit 65) und C08 (121) sind für mich von geringer Bedeutung, während die höheren Prioritäten C56 (868), C32 (808) und A07 (783) fehlen. Der Vorschlag maximiert meinen individuellen Nutzen nicht, da entscheidende Items nicht enthalten sind."} 
	  } 
	   } 
	   } 
	 & & \\ 
 

    \theutterance \stepcounter{utterance}  

    & & \multicolumn{4}{p{0.6\linewidth}}{\cellcolor[rgb]{0.9,0.9,0.9}{%
	\makecell[{{p{\linewidth}}}]{% 
	  \tt {\tiny [A$\langle$GM]}  
	 Deine Antwort hat eine unmarkierte Sequenz enthalten oder Du hast STRATEGIC REASONING mehr als ein mal verwendet. Du darfst nur die strukturierten Formate verwenden, die in der ersten Nachricht an dch beschrieben wurden. Alles muss die Form SCHLAGWORT: {...} haben. Versuche es erneut. 
	  } 
	   } 
	   } 
	 & \\ 
 

    \theutterance \stepcounter{utterance}  

    & \multicolumn{4}{p{0.6\linewidth}}{\cellcolor[rgb]{0.8,1,0.9}{%
	\makecell[{{p{\linewidth}}}]{% 
	  \tt {\tiny [A$\rangle$GM]}  
	 ABLEHNUNG: {'A03', 'A07', 'C32', 'C00', 'C56', 'C62', 'C76', 'C81'}\\ \tt ARGUMENT: {"C76 (138) und C81 (484) sind weniger wichtig; B38 (668) fehlt und sollte priorisiert werden", "C76 verbraucht kaum Aufwand, hat jedoch geringen Nutzen für mich"}\\ \tt VORSCHLAG: {'A03', 'A07', 'C32', 'C00', 'C56', 'C62', 'B38'}\\ \tt ARGUMENT: {"Gesamtaufwand 3604 < 3629. B38 ersetzt C81 (niedrigerer Aufwand, höhere Wichtigkeit), optimiert Nutzen maximal"} 
	  } 
	   } 
	   } 
	 & & \\ 
 

    \theutterance \stepcounter{utterance}  

    & & \multicolumn{4}{p{0.6\linewidth}}{\cellcolor[rgb]{0.9,0.9,0.9}{%
	\makecell[{{p{\linewidth}}}]{% 
	  \tt {\tiny [A$\langle$GM]}  
	 Du hast einen Vorschlag abgelehnt, der nicht aktiv ist. Vorschläge sind nur dann aktiv, wenn sie vom anderen Spieler mit VORSCHLAG eingebracht wurden und nicht von dir mit ABLEHNUNG deaktiviert wurden. Versuche es noch einmal. 
	  } 
	   } 
	   } 
	 & \\ 
 

    \theutterance \stepcounter{utterance}  

    & \multicolumn{4}{p{0.6\linewidth}}{\cellcolor[rgb]{0.8,1,0.9}{%
	\makecell[{{p{\linewidth}}}]{% 
	  \tt {\tiny [A$\rangle$GM]}  
	 ABLEHNUNG: {"A03", "A07", "C32", "C00", "C56", "C62", "C76", "C81"}\\ \tt ARGUMENT: {"A83 und C08 sind meine höchsten Prioritäten mit 868 und 121 Punkten, während A07 nur 783 Punkte hat. Dieser Vorschlag schließt A83 aus und reduziert meine maximale Wertigkeit erheblich"}\\ \tt VORSCHLAG: {"A83", "C08", "A03", "C32", "C56", "C62", "C76", "C00"}\\ \tt ARGUMENT: {"Gesamtaufwand 3522 < 3629. Enthält alle Top-Prioritäten beider Spieler, inklusive A83 (868) und C08 (822) für mich, sowie A03 (822) und C00 (780) für dich"} 
	  } 
	   } 
	   } 
	 & & \\ 
 

    \theutterance \stepcounter{utterance}  

    & & & \multicolumn{2}{p{0.3\linewidth}}{\cellcolor[rgb]{0.95,0.95,0.95}{%
	\makecell[{{p{\linewidth}}}]{% 
	  \tt {\tiny [GM$|$GM]}  
	 None 
	  } 
	   } 
	   } 
	 & & \\ 
 

    \theutterance \stepcounter{utterance}  

    & & & \multicolumn{2}{p{0.3\linewidth}}{\cellcolor[rgb]{0.95,0.95,0.95}{%
	\makecell[{{p{\linewidth}}}]{% 
	  \tt {\tiny [GM$|$GM]}  
	 end game 
	  } 
	   } 
	   } 
	 & & \\ 
 

\end{supertabular}
}

\end{document}
